% !TeX root = ../main.tex

\chapter{Praktijktesten op het circuit}
\label{chap:spa}

\section{Context van het eerste raceweekend}
Het dashboard werd tijdens een volledig raceweekend in Spa gebruikt met "RIS Timing" als timingprovider. Dit was de
eerste en enige officiële praktijktest tijdens een raceweekend vóór de \emph{24 Hours of Zolder}.
Vrijdag kon de
vrije training succesvol worden gevolgd en werd de ondersteuning voor een tweede timingprovider in
een echte omgeving getest. De kwalificatie verliep anders dan gepland: door motorschade kon de
voorziene wagen niet verder worden gebruikt. Het team bouwde gedurende de nacht een ouder Mazda
MX-5-model op om toch aan de race deel te nemen. Deze wagen had ongeveer 50 pk minder vermogen en
reed op oude banden uit noodzaak. De omstandigheden waren daardoor niet geschikt om de absolute prestaties van
de oorspronkelijk geplande wagen te beoordelen, maar leverden wel waardevolle circuitervaring voor
de nieuwe piloot en een realistische stresstest voor de applicatie.

Deze test was relevant met het oog op de 24 uur van Zolder, waarvoor het dashboard meerdere teams
moet kunnen volgen. Samen met Jan ondersteunde ik enkele strategische afwegingen. Daarnaast gaf hij
gerichte feedback over de gebruikersinterface, die ik na het raceweekend meteen kon verwerken.

\section{Uitgevoerde tests en verzamelde data}
Tijdens practice, pre-qualifying, qualifying en race werden afzonderlijke sessiemappen bijgehouden. De opgeslagen data bestond uit rondehistoriek, analytics,
rondevoorspellingen, pitstopplannen en sessiemetadata. Tijdens de race werden onder andere de
volgende onderdelen gebruikt:

\begin{itemize}
  \item parsing van RIS Timing, inclusief sessiegegevens, pilootnamen, sectoren, GAP en INT
  \item langdurige opslag en hervatting van rondehistoriek
  \item piloot-, team- en klassevergelijkingen
  \item rondevoorspelling en referentietijden
  \item gaps naar klassewagens en catchindicaties
  \item verplichte pitstops en de voorspelde positie na een pitstop
  \item grafieken en meerdere gevolgde wagens
\end{itemize}

De test bewees dat de volledige keten gedurende een echte sessie data kon verzamelen en bewaren. Ze
bracht tegelijk fouten aan het licht die niet voldoende zichtbaar waren in demo's en zelfgeschreven unit tests. De Spa-data wordt daarom niet alleen als resultaat bewaard, maar ook als regressiedataset voor volgende versies.

\section{Belangrijkste bevindingen}
\subsection{Vlagcontext en geldige rondes}
De belangrijkste datakwaliteitsfout was het moment waarop FCY-status aan een ronde
en sector wordt gekoppeld. Wanneer FCY begint terwijl een wagen in sector 2 rijdt, moet sector 2
onmiddellijk als geneutraliseerd worden gemarkeerd. Sector 1 blijft geldig wanneer die volledig
onder groen werd gereden maar de volledige ronde mag niet meer meetellen als representatieve rondetijd.

Een concreet voorbeeld kwam in ronde 2. In \code{lap_history.jsonl} staat de afgeronde ronde
met \code{lapFlag = Green flag}, terwijl de opgeslagen sectorflags FCY aangaven. De sectoren werden
hierdoor al uitgesloten, maar de rondecontext is intern niet consistent. Daarnaast moeten pit-in- en
outlaps en duidelijke tempo-outliers door exact dezelfde centrale selectieregel worden uitgesloten.
De test toonde dus dat een globale sessieflag op het moment van afronden niet volstaat maar dat de applicatie
ook de sector waarin iedere gevolgde wagen zich bevindt moet onthouden wanneer de vlag verandert.

\subsection{Gemiddelden en vergelijkingseenheden}
De getoonde gemiddelden waren op enkele momenten hoger dan een handmatige controle deed verwachten.
Waarschijnlijke oorzaken zijn onvolledige vlagcontext en pitgerelateerde ronden die niet overal
identiek gefilterd werden. Bovendien kon dezelfde wagen tegelijk als beste wagen in de klasse en als
geselecteerde vergelijkingswagen verschijnen, terwijl beide vakken toch een ander gemiddelde
toonden. Dat verschil ontstaat wanneer het ene vak een volledige wagenhistoriek gebruikt en het
andere een actuele piloot of stint, zonder dit duidelijk te benoemen.
De gewenste oplossing is één volledig analyticsdocument met benoemde statistieken en
selectiecriteria. Dashboard, grafieken en rapportage moeten dezelfde waarden uit dat document
gebruiken.

\subsection{Gaps en pitstopprojectie}
Gaps naar de voorliggende en achterliggende klassewagen veranderden tijdens de race te sterk om
operationeel betrouwbaar te zijn. RIS Timing werkte met GAP tot de leider en INT tot de vorige
wagen dus tussenliggende wagens uit andere klassen moesten in de intervalketen worden opgenomen.
Rondeachterstanden maakten een directe omzetting bovendien grof en gevoelig voor het meetmoment.
Voor volgende iteraties is een gapgeheugen nodig. Een gap wordt pas als nieuwe stabiele waarde
gepubliceerd wanneer de betrokken wagen een nieuw meetpunt passeert, bij voorkeur de
start-finishlijn. De UI moet groot de laatst bevestigde gap tonen, met daaronder het verschil tussen
de recentste rondetijden en een afzonderlijke catchtrend. Een wagen die al enige tijd (5 van onze rondes hebben we voor nu gekozen) in de pits staat, moet tijdelijk als niet-actieve tegenstander worden gemarkeerd in
plaats van een misleidende catchtijd te krijgen. Dezelfde stabiele gapreeks kan later de
pitstopprojectie en een gap-over-timegrafiek voeden.


\section{Nieuwe requirements uit de praktijktest}
De test leverde naast correcties ook drie nieuwe uitbreidingswensen op:

\begin{enumerate}
  \item Een piloottimer met duur van de huidige stint en totale rijtijd per piloot. Want ook dit is belangrijk tijdens endurance races, piloten mogen niet langer dan een bepaalde periode in de auto zitten volgens de regels. Pilootwissels
        vormen hier de stintgrens, sessietijd en timestamps leveren de duur.
  \item Een scrollbare recente-rondenstrook met rondenummer en kleurstatus: pit-in met rode
        markering en letter P, outlap rood, FCY/SC geel, snelste eigenronde groen en als deze de snelste ronde van alle auto's in de klasse is deze paars (dit zijn conventies uit alle niveaus in de motorsport). De volledige historie blijft beschikbaar,
        terwijl standaard slechts de laatste tien tot twintig ronden zichtbaar zijn.
  \item Automatische PDF-rapportage per stint en na de race. Een stintdocument bevat ronde- en
        sectortijden, geldige gemiddelden, vergelijking met teamgenoten en de gapevolutie. Bij een
        pilootwissel kan automatisch een map zoals bijvoorbeeld \code{STINT_2_JANSSENS_Robbe} worden gemaakt. Een
        raceoverzicht combineert de stintresultaten per piloot en voor het team.
\end{enumerate}

Deze rapportage is technisch haalbaar omdat de onderliggende gegevens al lokaal worden opgeslagen.
Voor betrouwbare automatische generatie moeten eerst stintdetectie, vlagcontext en centrale
analytics eenduidig zijn, anders zou een verzorgd PDF-document dezelfde foutieve selectie
reproduceren.

\section{Prioriteiten na Spa}
De vervolgvolgorde wordt bepaald door databetrouwbaarheid, niet door zichtbaarheid in de interface:

\begin{enumerate}
  \item vlagtransities per wagen en sector correct registreren en alle pacefilters centraliseren
  \item één persistent analyticsmodel invoeren
  \item stabiele gaps op meetpunten bewaren en gebruiken voor catch- en pitstopprojecties
  \item stintdetectie, pilootduur en recente-rondenstrook toevoegen
  \item automatische stint- en racerapporten genereren
  \item pas daarna de geplande UI-herwerking uitvoeren op basis van dezelfde opgeslagen
        viewmodellen
\end{enumerate}

\section{Aanvullende tests tussen Spa en Zolder}
Omdat Spa de enige officiële praktijktest was, konden nieuwe functies die daarna werden toegevoegd
niet meer tijdens een vergelijkbaar raceweekend worden getest. Dat maakte het moeilijk om met
zekerheid te bepalen of nieuwe berekeningen en interfacewijzigingen ook onder echte racedruk correct
zouden blijven werken. De opgeslagen Spa-data bleef bruikbaar voor regressietests, maar bevatte
uiteraard niet elke situatie die tijdens een 24-uursrace kan voorkomen.
Om toch voortdurend met veranderende timingdata te testen, gebruikte ik daarom veelvuldig de live
demo van GetRaceResults\footnote{https://livetiming.getraceresults.com/demo\#screen-results}. Die demo maakte het mogelijk om de parser, de live updates en de
dashboardweergave herhaaldelijk te controleren zonder op een volgend officieel raceweekend te moeten
wachten. Ze vervangde echter geen volledige praktijktest want de demo bood uiteraard niet dezelfde combinatie van
meerdere teams, pitstops, neutralisaties, pilootwissels en operationele tijdsdruk als tijdens een
echte endurancewedstrijd.
Daarnaast werden gerichte unit tests gebruikt om zoveel mogelijk edge cases af te dekken. Daarbij
werd onder meer gecontroleerd hoe de applicatie omgaat met ontbrekende of onvolledige sectoren,
pit-in- en outlaps, vlagovergangen, rondeachterstanden, wisselende piloten en tegenstanders die
langdurig in de pits staan. Deze tests vergrootten het vertrouwen in afzonderlijke berekeningen,
maar konden niet volledig aantonen hoe alle componenten gedurende 24 uur samen zouden functioneren.
Juist daardoor bleef er voor de \emph{24 Hours of Zolder} de nodige spanning want het was tegelijk het
einddoel van de stage en de eerste keer dat de nieuwste functies in hun volledige gebruikscontext
zouden worden ingezet.

% !TeX root = ../main.tex

\section{Het einddoel: \emph{24 Hours of Zolder}}
\label{chap:24hofZolder}

\pending{Na de \emph{24 Hours of Zolder}: vervang het voorlopige racescenario door de werkelijke gebeurtenissen en resultaten. Beschrijf ook welke functies tijdens de race zijn gebruikt, welke beslissingen het dashboard ondersteunde, welke fouten of beperkingen zichtbaar werden en welke feedback Jan en de piloten gaven.}

\subsection{Verwachte meerwaarde tijdens de race}
De grootste meerwaarde tijdens de 24 uur van Zolder lag in het combineren van informatie die op de
timingpagina verspreid of impliciet aanwezig is. Het dashboard toont niet alleen de actuele positie,
maar vertaalt die naar teamgerichte vragen:

\begin{itemize}
  \item Hoe snel rijdt de huidige piloot tegenover zijn teamgenoten?
  \item Hoe evolueert het tempo tegenover de beste wagen in de klasse?
  \item Hoeveel marge is er tegenover de referentietijd?
  \item Wanneer wordt een pitstop reglementair mogelijk of net riskant?
  \item Waar komt de wagen vermoedelijk terug op de baan na een pitstop?
  \item Welke rondes of sectoren zijn representatief genoeg om in gemiddelden te gebruiken?
\end{itemize}

Deze vragen zijn vooral belangrijk omdat Jan tijdens de race meerdere rollen tegelijk opneemt:
strategie, coaching, communicatie en opvolging van het wedstrijdverloop. Een dashboard dat
herhalende berekeningen automatisch uitvoert, verkleint de kans dat belangrijke details worden
gemist wanneer er veel tegelijk gebeurt. De applicatie neemt geen beslissingen in de plaats van de
race-engineer, maar maakt de beschikbare informatie sneller interpreteerbaar.

\subsection{Wat tijdens de 24 uur getest moest worden}
Hoewel de applicatie voor de race klaar was als werkende basis, moest de 24 uur van Zolder nog
aantonen hoe goed ze standhield in haar uiteindelijke context. De belangrijkste validatiepunten
waren:

\begin{itemize}
  \item Langdurige stabiliteit gedurende een volledige 24-uursrace.
  \item Correcte opvolging van meerdere wagens tegelijk.
  \item Betrouwbare stintdetectie bij alle pilootwissels.
  \item Correcte behandeling van FCY, Safety Car, rode vlag en pitgerelateerde rondes.
  \item Bruikbaarheid van de pitstopplanner wanneer snel beslist moet worden.
  \item Leesbaarheid van het dashboard in een drukke raceomgeving.
  \item Nut van de automatische stint- en racerapporten na de wedstrijd.
\end{itemize}

Vooral de combinatie van een lange looptijd, meerdere pitstops, wisselende wedstrijdsituaties en
menselijke werkdruk kon niet volledig met korte tests worden nagebootst. De 24 uur van Zolder was
daarom niet alleen het einddoel van de applicatie, maar ook de belangrijkste test van de
ontwerpkeuzes.

\begingroup
\color{pendingred}
\subsection{Voorlopig racescenario -- na de race te verifiëren}

\textbf{De onderstaande tekst is een voorlopig scenario. Alle resultaten, gebeurtenissen en
conclusies moeten na de race worden gecontroleerd en vervangen of bevestigd met de werkelijke
gegevens.}
\pending{TODO!!!}

In dit voorlopige scenario startte de wagen vanaf poleposition. Enkele fouten tijdens de eerste uren
kostten het team echter posities, waardoor de race al vroeg minder vanzelfsprekend verliep dan de
startpositie deed vermoeden. De analyses in de applicatie bleven gedurende die fase een duidelijk
overzicht geven van het eigen tempo, de relevante klassewagens en de verwachte gevolgen van
strategische keuzes.

Op een belangrijk moment liet de pitstopprojectie zien dat een pitstop een gunstige keuze was. De
verwachting was dat de wagen na de stop net vóór de achterliggende concurrent opnieuw op het circuit
zou komen. Daardoor kon het team de beslissing niet alleen op gevoel nemen, maar ook ondersteunen
met de actuele gaps, rondetijden en de berekende rejoinpositie uit het dashboard.

Het tempo lag uiteindelijk lager dan dat van enkele andere wagens en eerdere fouten bleven invloed
hebben op het resultaat. Door het uitvallen van een concurrent eindigde het team volgens dit
scenario toch op de derde plaats. Vanuit poleposition was dat niet het oorspronkelijk gehoopte
resultaat, maar gezien het wedstrijdverloop kon het team tevreden terugkijken op een podiumplaats.

Het raceweekend was bovendien belangrijk voor de zichtbaarheid van MSTC en Mazda. Het BV-team, mede
onder leiding van Bert Longin, zorgde voor extra publieke aandacht. Daardoor bood de wedstrijd niet
alleen sportieve en technische validatie voor de applicatie, maar ook publicitaire meerwaarde voor
de betrokken organisaties.
\endgroup

